\documentclass[11pt,a4paper]{article}
\usepackage[utf-8]{inputenc}
\usepackage[spanish]{babel}
\usepackage{geometry}
\usepackage{graphicx}
\usepackage{hyperref}
\usepackage{listings}
\usepackage{xcolor}
\usepackage{fancyhdr}
\usepackage{float}

% Márgenes
\geometry{left=2cm, right=2cm, top=2cm, bottom=2cm}

% Colores para código
\definecolor{codegreen}{rgb}{0,0.6,0}
\definecolor{codegray}{rgb}{0.5,0.5,0.5}
\definecolor{codepurple}{rgb}{0.58,0,0.82}
\definecolor{backcolour}{rgb}{0.95,0.95,0.92}

% Configuración de listings para código Dart
\lstdefinestyle{dartStyle}{
    backgroundcolor=\color{backcolour},
    commentstyle=\color{codegreen},
    keywordstyle=\color{codepurple},
    numberstyle=\tiny\color{codegray},
    stringstyle=\color{codepurple},
    basicstyle=\ttfamily\small,
    breakatwhitespace=false,
    breaklines=true,
    captionpos=b,
    keepspaces=true,
    numbers=left,
    numbersep=5pt,
    showspaces=false,
    showstringspaces=false,
    showtabs=false,
    tabsize=2
}

\lstset{style=dartStyle}

% Encabezados y pies
\pagestyle{fancy}
\fancyhf{}
\rhead{Servicios Básicos}
\cfoot{\thepage}

% Título
\title{\textbf{Aplicación de Pagos de Servicios Básicos} \\ \textit{Flutter - Tarea 1.2}}
\author{}
\date{}

\begin{document}

\maketitle
\vspace{-1cm}

% Sección 1: Introducción
\section{Introducción}

La presente aplicación móvil es una solución desarrollada en \textbf{Flutter} que permite registrar y calcular pagos de servicios básicos frecuentes. La aplicación implementa patrones de arquitectura moderna como \textbf{MVC (Modelo-Vista-Controlador)} y \textbf{Atomic Design}, garantizando un código escalable, mantenible y reutilizable.

La aplicación está dirigida a usuarios que deseen:
\begin{itemize}
    \item Registrar información de clientes
    \item Seleccionar tipos de servicios (agua, energía, internet, TV)
    \item Calcular montos a pagar basados en consumo y tarifas
    \item Aplicar descuentos o recargos
    \item Obtener comprobantes detallados de pagos
\end{itemize}

\section{Descripción del Problema}

\subsection{Objetivo General}

Crear una aplicación móvil que automatice el cálculo de pagos de servicios básicos, permitiendo a los usuarios registrar datos de clientes, seleccionar servicios, ingresar consumos y obtener un resumen detallado del pago con descuentos y recargos aplicables.

\subsection{Servicios Considerados}

La aplicación soporta los siguientes servicios:

\begin{table}[H]
\centering
\begin{tabular}{|l|c|l|}
\hline
\textbf{Servicio} & \textbf{Tarifa Base} & \textbf{Unidad} \\
\hline
Agua Potable & \$0.50 & por m³ \\
\hline
Energía Eléctrica & \$0.15 & por kWh \\
\hline
Internet y Telefonía & \$25.00 & valor fijo \\
\hline
TV por Cable & \$20.00 & valor fijo \\
\hline
\section{Descripción del Problema}

La aplicación calcula pagos de servicios básicos: agua (\$0.50/m³), energía (\$0.15/kWh), internet y telefonía (\$25), TV por cable (\$20). 

\textbf{Funcionalidades:} Registrar cliente (nombre, cédula, dirección), seleccionar servicio, ingresar consumo, aplicar descuentos/recargos y generar comprobante con cálculo: \textit{Total} = (Consumo × Tarifa) - Descuento + Recargo.

\vspace{0.5cm}
\noindent\fbox{\begin{minipage}{\textwidth}
\textbf{Captura recomendada 1:} Pantalla principal del programa ejecutándose con el formulario completo visible.
\end{minipage}t                    # Punto de entrada
├── model/
│   └── servicio_model.dart     # Modelos de datos
├── controller/
│   └── servicio_controller.dart # Lógica de negocio
├── view/
│   ├── home_servicios_view.dart           # Pantalla principal
│   └── resultado_servicios_view.dart      # Pantalla de resultados
└── widgets/
    ├── atoms/                   # Componentes básicos
    │   ├── input_text.dart
    │   ├── custom_label.dart
    │   ├── custom_button.dart
    │   ├── custom_checkbox.dart
    │   └── custom_radio.dart
    ├── molecules/              # Combinaciones de átomos
    │   ├── info_cliente.dart
    │   ├── datos_consumo.dart
    │   ├── selector_servicio.dart
    │   ├── opciones_descuento.dart
    │   └── opciones_recargo.dart
    └── organisms/              # Combinaciones de moléculas
        └── formulario_pago.dart
\end{lstlisting}

\section{Componentes del Sistema}

\subsection{Modelos de Datos}

Los modelos definen la estructura de los datos utilizados en la aplicación:

\subsubsection{ServicioModel}

Define la estructura básica de un servicio:

\begin{lstlisting}[language=Dart]
class ServicioModel {
  final String nombre;
  final double tarifaBase;
  final double unidad;
  
  ServicioModel({
    required this.nombre,
    required this.tarifaBase,
    required this.unidad,
  });
}
\end{lstlisting}
MVC + Atomic Design}

\subsection{Patrón MVC}

La aplicación separa responsabilidades en:

\begin{itemize}
    \item \textbf{Model:} \texttt{servicio\_model.dart} - Define \texttt{ClienteModel}, \texttt{ServicioModel}, \texttt{PagoServicioModel}
    \item \textbf{Controller:} \texttt{servicio\_controller.dart} - Valida datos y calcula pagos
    \item \textbf{View:} \texttt{home\_servicios\_view.dart}, \texttt{resultado\_servicios\_view.dart} - Interfaz gráfica
\end{itemize}

\subsection{Atomic Design}

Estructura de carpetas:
\begin{lstlisting}[basicstyle=\ttfamily\tiny]
lib/
├── main.dart
├── model/ → servicio_model.dart
├── controller/ → servicio_controller.dart
├── view/ → home_servicios_view.dart
│          resultado_servicios_view.dart
└── widgets/
    ├── atontrolador y Lógica de Negocio}

El \texttt{ServicioController} implementa:

\begin{itemize}
    \item \textbf{Tarifas}: Mapa con precios de servicios
    \item \textbf{Validaciones}: Nombre, cédula, dirección, consumo positivo
    \item \textbf{Cálculo}: Subtotal = Consumo × Tarifa
    \item \textbf{Descuento/Recargo}: Porcentaje sobre subtotal
    \item \textbf{Total}: Subtotal - Descuento + Recargo
\end{itemize}

Ejemplo de cálculo: Si agua = 50m³ × \$0.50 = \$25.00, con descuento 10\% y recargo 5\%:
\begin{center}
Total = \$25.00 - \$2.50 + \$1.25 = \$23.75
\end{center}

\vspace{0.3cm}
\noindent\fbox{\begin{minipage}{\textwidth}
\textbf{Captura recomendada 3:} Código del método \texttt{procesarPago()} en el controlador (líneas de cálculo principal).
\end{minipage}
\end{tabular}
\caption{Componentes moleculares}
\end{table}

\subsection{Organismos}

Composiciones complejas de moléculas:

\begin{itemize}
    \item \textbf{FormularioPago:} Combina todas las moléculas en un formulario completo
\end{itemize}

\section{Flujo de Funcionamiento}

\subsection{Diagrama de Flujo}

\begin{enumerate}
    \item Usuario abre la aplicación (HomeServiciosView)
    \item Ingresa datos del cliente (nombre, cédula, dirección)
    \item Selecciona tipo de servicio
    \item Ingresa consumo o valor
    \item Opcionalmente aplica descuento o recargo
    \item Presiona botón "Calcular"
    \item Controlador valida y procesa la información
    \item Se genera comprobante
    \item Navega a ResultadoServiciosView
    \item Muestra comprobante al usuario
    \item Usuario puede volver atrás para nuevo registro
\end{enumerate}

\subsection{Ejemplo de Cálculo}

Supongamos el siguiente escenario:

\begin{table}[H]
\centering
\begin{tabular}{|l|c|}
\hline
\textbf{Concepto} & \textbf{Valor} \\
\hline
Nombre Cliente & Juan García \\
\hline
Servicio & Agua Potable \\
\hline
Consumo & 50 m³ \\
\hline
Tarifa & \$0.50 por m³ \\
\hline
\end{tabular}
\caption{Datos de entrada}
\end{table}

\textbf{Cálculos:}
y Widgets}

\subsection{Vistas (Views)}

\begin{itemize}
    \item \textbf{HomeServiciosView}: Formulario principal con campos de cliente, servicio, consumo, descuento/recargo
    \item \textbf{ResultadoServiciosView}: Muestra comprobante detallado con todos los cálculos
\end{itemize}

\subsection{Atomic Design - Componentes}

\textbf{Átomos:} InputText, CustomLabel, CustomButton, CustomCheckbox, CustomRadio

\textbf{Moléculas:} InfoCliente (3×InputText), SelectorServicio (5×CustomRadio), OpcionesDescuento (CustomCheckbox+InputText)

\textbf{Organismos:} FormularioPago (combinación de moléculas)

\vspace{0.3cm}
\noindent\fbox{\begin{minipage}{\textwidth}
\textbf{Captura recomendada 4:} Pantalla del programa con formulario lleno de datos listos para calcular.
\end{minipage}scuentos (reducen el monto)
    \item Recargos (aumentan el monto)
\end{itemize}

\section{Tecnologías Utilizadas}

\begin{itemize}
    \item \textbf{Flutter:} Framework para desarrollo multiplataforma
    \item \textbf{Dart:} Lenguaje de programación
    \item \textbf{Material Design:} Sistema de diseño Google
    \item \textbf{Navigator:} Sistema de rutas de Flutter
\end{itemize}

\section{Requisitos Técnicos Cumplidos}

\begin{table}[H]
\centering
\begin{tabular}{|l|c|}
\hline
\textbf{Requisito} & \textbf{Estado} \\
\hline
Patrón MVC & ✓ Implementado \\
\hline
Atomic Design & ✓ Implementado \\
\hline
Sistema de Rutas & ✓ Implementado \\
\hline
Separación de Carpetas & ✓ Implementado \\
\hline
Cajas de Texto & ✓ Implementadas \\
\hline
Labels & ✓ Implementados \\
\hline
Botones & ✓ Implementados \\
\hline
CheckBox & ✓ Implementados \\
\hline
RadioButton & ✓ Implementados \\
\hline
Interfaz Clara y Funcional & ✓ Implementada \\
\hline
\end{tabular}
\caption{Estado de requisitos}
\end{table}

\section{Conclusiones}

La aplicación de \textbf{Pagos de Servicios Básicos} representa una solución completa y escalable para el cálculo automatizado de tarifas de servicios. 

\subsection{Logros Principales}

\begin{itemize}
    \item Implementación exitosa del patrón MVC
    \item Aplicación consistente de Atomic Design
    \item Interfaz intuitiva y funcional
    \item Lógica de negocio clara y validable
    \item Código modular y reutilizable
    \item Sistema de navegación robusto
\end{itemize}

\subsection{Ventajas del Diseño}

\begin{itemize}
    \item \texty Funcionamiento}

\begin{enumerate}
    \item Usuario ingresa datos (nombre, cédula, dirección)
    \item Selecciona servicio con RadioButton
    \item Ingresa consumo/valor y opcionales descuento/recargo
    \item Presiona "Calcular"
    \item Controller valida y genera comprobante
    \item Navega a pantalla de resultados
\end{enumerate}

\vspace{0.3cm}
\noindent\fbox{\begin{minipage}{\textwidth}
\textbf{Captura recomendada 5:} Pantalla de resultados mostrando el comprobante completo generado.
\end{minipage}}

\vspace{0.5cm}

\section{Requisitos Cumplidos}

\begin{itemize}
    \item ✓ MVC: Modelos, Controlador, Vistas separados
    \item ✓ Atomic Design: Átomos, Moléculas, Organismos
    \item ✓ Rutas: push, pop, pushNamed
    \item ✓ Cajas de texto, Labels, Botones, CheckBox, RadioButton
    \item ✓ Validaciones de datos
    \item ✓ Cálculos automáticos
    \item ✓ Generación de comprobante